% math4cs-hw3-predicate-logic-solution.tex

% !TEX program = xelatex
%%%%%%%%%%%%%%%%%%%%
% see http://mirrors.concertpass.com/tex-archive/macros/latex/contrib/tufte-latex/sample-handout.pdf
% for how to use tufte-handout
\documentclass[a4paper, justified]{tufte-handout}

\input{hw-preamble} % feel free to modify this file if you understand LaTeX well
%%%%%%%%%%%%%%%%%%%%
\title{计算机数学 (2-predicate-logic: 一阶谓词逻辑 II)}
\me{魏恒峰}{hfwei@hnu.edu.cn}{}{}
\date{2026年3月30日}
%%%%%%%%%%%%%%%%%%%%
\begin{document}
\maketitle
%%%%%%%%%%%%%%%%%%%%
\noplagiarism % PLEASE DON'T DELETE THIS LINE!
%%%%%%%%%%%%%%%%%%%%
\begin{abstract}
  % \mfigcap{width = 0.85\textwidth}{figs/George-Boole}{George Boole}
  % \begin{center}{\fcolorbox{blue}{yellow!60}{\parbox{0.65\textwidth}{\large
  %   \begin{itemize}
  %     \item
  %   \end{itemize}}}}
  % \end{center}
\end{abstract}
%%%%%%%%%%%%%%%%%%%%
\beginrequired
%%%%%%%%%%%%%%%

%%%%%%%%%%%%%%%
\newpage
\begin{problem}[一阶谓词逻辑: 模型与语义]
  考虑一阶谓词逻辑逻辑的闭公式 (也称为语句; sentence)：
  \[
    \phi \triangleq \forall x\; \exists y\; \exists z\;
      (P(x, y) \land P(z, y) \land (P(x, z) \to P(z, x)))
  \]
  以下哪些模型 $\mathcal{M}$ 满足 $\phi$? 请给出理由。
  \begin{enumerate}[(1)]
    \item $\mathcal{M}_{1}$:
      $|\mathcal{M}| \triangleq \mathbb{N}$,
      $P^{\mathcal{M}} \triangleq \{(x, y) \mid x < y\}$.
    \item $\mathcal{M}_{2}$:
      $|\mathcal{M}| \triangleq \mathbb{N}$,
      $P^{\mathcal{M}} \triangleq \{(x, y) \mid x < y + 1\}$.
  \end{enumerate}
\end{problem}

\begin{solution}
  \begin{intention}
    逐个模型为任意的 $x$ 构造见证元 $y, z$; 关键是注意第三个合取项是蕴含式, 可以通过让前件为假来满足。
  \end{intention}

  由于 $\phi$ 是闭公式, 判定 $\mathcal{M} \models \phi$ 与变元赋值无关。

  \begin{enumerate}[(1)]
    \item 对 $\mathcal{M}_{1}$, 有
      \[
        P^{\mathcal{M}_{1}}(a, b) \iff a < b.
      \]
      任取 $x \in \mathbb{N}$, 取
      \[
        y \triangleq x + 1,
        \qquad
        z \triangleq x.
      \]
      则
      \[
        P(x, y) \iff x < x + 1,
        \qquad
        P(z, y) \iff x < x + 1,
      \]
      均成立; 且
      \[
        P(x, z) \to P(z, x)
        \iff
        (x < x) \to (x < x)
      \]
      也成立, 因为前件 $x < x$ 为假。
      故对任意 $x$ 都存在这样的 $y, z$, 从而
      \[
        \mathcal{M}_{1} \models \phi.
      \]

      \textbf{分值分配（10分）.}
      \begin{enumerate}[(a)]
        \item 正确写出 $P(a, b) \iff a < b$: \score{2}
        \item 为任意 $x$ 选出见证元 $y = x + 1$: \score{2}
        \item 为任意 $x$ 选出见证元 $z = x$: \score{2}
        \item 验证三个合取项都成立: \score{3}
        \item 得出 $\mathcal{M}_{1} \models \phi$: \score{1}
      \end{enumerate}

    \item 对 $\mathcal{M}_{2}$, 有
      \[
        P^{\mathcal{M}_{2}}(a, b) \iff a < b + 1,
      \]
      亦即在 $\mathbb{N}$ 上等价于 $a \le b$。
      任取 $x \in \mathbb{N}$, 取
      \[
        y \triangleq x,
        \qquad
        z \triangleq x.
      \]
      则
      \[
        P(x, y) \iff x \le x,
        \qquad
        P(z, y) \iff x \le x,
      \]
      均成立; 且
      \[
        P(x, z) \to P(z, x)
        \iff
        (x \le x) \to (x \le x)
      \]
      显然成立。
      故对任意 $x$ 都存在这样的 $y, z$, 从而
      \[
        \mathcal{M}_{2} \models \phi.
      \]

      \textbf{分值分配（10分）.}
      \begin{enumerate}[(a)]
        \item 正确把 $a < b + 1$ 化为 $a \le b$: \score{2}
        \item 为任意 $x$ 选出见证元 $y = x$: \score{2}
        \item 为任意 $x$ 选出见证元 $z = x$: \score{2}
        \item 验证三个合取项都成立: \score{3}
        \item 得出 $\mathcal{M}_{2} \models \phi$: \score{1}
      \end{enumerate}
  \end{enumerate}

  \begin{remark}
    \textbf{重点.} 把“满足闭公式”转化为“对任意 $x$ 能构造见证元”。
    \textbf{难点.} 第三个条件是蕴含式, 并不要求 $P(x, z)$ 真; 让前件为假常常更省事。
    \textbf{易错点.} 在 $\mathcal{M}_{2}$ 中, $x < y + 1$ 在 $\mathbb{N}$ 上等价于 $x \le y$, 不是严格小于。
  \end{remark}
\end{solution}
%%%%%%%%%%%%%%%

%%%%%%%%%%%%%%%
\newpage
\begin{problem}[一阶谓词逻辑: 形式化描述与推理]
  请使用一阶谓词逻辑公式描述以下两个定义,
  并从逻辑推理的角度说明这两种定义之间是否有强弱之分
  （可以使用``重言式''进行公式转换）。
  \noindent {(要求: 需给出关键的推理步骤或理由)}

  A function $f$ from $\mathbb{R}$ to $\mathbb{R}$ is called
  \begin{enumerate}[(1)]
    \setlength{\itemsep}{10pt}
    \item \blue{\it pointwise continuous} (连续的) if
      for every $x \in \mathbb{R}$
      and every real number $\epsilon > 0$,
      there exists real $\delta > 0$ such that
      for every $y \in \mathbb{R}$ with $|x - y| < \delta$,
      we have that $|f(x) -  f(y)|< \epsilon$.
    \item \blue{\it uniformly continuous} (一致连续的) if
      for every real number $\epsilon > 0$,
      there exists real $\delta > 0$ such that
      for every $x, y \in \mathbb{R}$ with $|x - y| < \delta$,
      we have that $|f(x) -  f(y)|< \epsilon$.
  \end{enumerate}
\end{problem}

\begin{solution}
  \begin{intention}
    先把“$\delta$ 是否依赖于 $x$”写成量词次序的差异, 再由量词前束判断两种定义的强弱关系。
  \end{intention}

  取论域为 $\mathbb{R}$。
  语言中包含常元 $0$, 函数符号 $f$, 减法 $-$, 绝对值 $|\cdot|$, 以及二元关系符号 $<$。

  \begin{enumerate}[(1)]
    \item 点态连续（pointwise continuous）可形式化为
      \[
        \mathrm{PC}(f) \triangleq
        \forall x\; \forall \epsilon\; \Bigl(
          0 < \epsilon \to
          \exists \delta\; \bigl(
            0 < \delta
            \land
            \forall y\; (|x - y| < \delta \to |f(x) - f(y)| < \epsilon)
          \bigr)
        \Bigr).
      \]

    \item 一致连续（uniformly continuous）可形式化为
      \[
        \mathrm{UC}(f) \triangleq
        \forall \epsilon\; \Bigl(
          0 < \epsilon \to
          \exists \delta\; \bigl(
            0 < \delta
            \land
            \forall x\; \forall y\; (|x - y| < \delta \to |f(x) - f(y)| < \epsilon)
          \bigr)
        \Bigr).
      \]
  \end{enumerate}

  为了比较二者强弱, 记
  \[
    A(x, \epsilon, \delta) \triangleq
    \forall y\; (|x - y| < \delta \to |f(x) - f(y)| < \epsilon).
  \]
  则
  \[
    \mathrm{PC}(f)
    \equiv
    \forall x\; \forall \epsilon\; (0 < \epsilon \to \exists \delta\; (0 < \delta \land A(x, \epsilon, \delta))),
  \]
  \[
    \mathrm{UC}(f)
    \equiv
    \forall \epsilon\; (0 < \epsilon \to \exists \delta\; (0 < \delta \land \forall x\; A(x, \epsilon, \delta))).
  \]

  下面说明
  \[
    \mathrm{UC}(f) \vdash \mathrm{PC}(f).
  \]
  任取 $x \in \mathbb{R}$ 与 $\epsilon > 0$。
  由 $\mathrm{UC}(f)$, 可得某个 $\delta > 0$ 使得
  \[
    \forall u\; A(u, \epsilon, \delta).
  \]
  对该全称式作 $\forall$-消去, 取 $u := x$, 得
  \[
    A(x, \epsilon, \delta),
  \]
  即
  \[
    \forall y\; (|x - y| < \delta \to |f(x) - f(y)| < \epsilon).
  \]
  因而同一个 $\delta$ 就可作为点 $x$ 处关于 $\epsilon$ 的见证元。
  既然 $x, \epsilon$ 任意, 故
  \[
    \mathrm{UC}(f) \vdash \mathrm{PC}(f).
  \]

  反之,
  \[
    \mathrm{PC}(f) \nvdash \mathrm{UC}(f).
  \]
  根本原因是: 在 $\mathrm{PC}(f)$ 中, 量词前束为
  \[
    \forall x\; \forall \epsilon\; \exists \delta\; \cdots,
  \]
  这里的 $\delta$ 可以依赖于 $x$;
  而在 $\mathrm{UC}(f)$ 中, 前束为
  \[
    \forall \epsilon\; \exists \delta\; \forall x\; \cdots,
  \]
  这里要求同一个 $\delta$ 同时适用于所有 $x$。
  这显然更强。

  语义上, 反向蕴含也确实不成立。例如函数
  \[
    f(x) = x^{2}
  \]
  在 $\mathbb{R}$ 上点态连续, 但并非一致连续。

  综上,
  \[
    \mathrm{UC}(f) \text{ 比 } \mathrm{PC}(f) \text{ 更强},
    \qquad
    \mathrm{PC}(f) \text{ 比 } \mathrm{UC}(f) \text{ 更弱}.
  \]

  \begin{remark}
    \textbf{重点.} 两个定义的本质差异只在量词次序: $\delta$ 是否允许依赖于 $x$。
    \textbf{难点.} 不能把 $\forall x\, \exists \delta$ 与 $\exists \delta\, \forall x$ 互换。
    \textbf{易错点.} 一致连续要求“同一个 $\delta$ 对所有 $x$ 都有效”, 这比普通连续强得多。
  \end{remark}

  \textbf{分值分配（10分）.}
  \begin{enumerate}[(a)]
    \item 正确写出点态连续的量词结构: \score{3}
    \item 正确写出一致连续的量词结构: \score{3}
    \item 说明 $\mathrm{UC}(f) \vdash \mathrm{PC}(f)$: \score{2}
    \item 说明反向不成立的量词原因: \score{1}
    \item 给出语义反例或明确结论: \score{1}
  \end{enumerate}
\end{solution}
%%%%%%%%%%%%%%%

%%%%%%%%%%%%%%%
\newpage
\begin{problem}[一阶谓词逻辑: 自然推理系统]
  给定如下``前提'', 请判断``结论''是否有效。
  请在一阶谓词逻辑的自然推理系统中完成。

  {\bf 前提:}
  \begin{enumerate}[(1)]
    \item 每个人或者喜欢美剧, 或者喜欢韩剧 (可以同时喜欢二者);
    \item 任何人如果他喜欢抗日神剧, 他就不喜欢美剧;
    \item 有的人不喜欢韩剧。
  \end{enumerate}

  \vspace{0.20cm}
  \noindent {\bf 结论:} 有的人不喜欢抗日神剧 ({\it 幸亏如此})。
\end{problem}

\begin{solution}
  \begin{intention}
    先把三条前提形式化; 然后从“不喜欢韩剧”的见证人出发, 用析取消去和反证法推出他不喜欢抗日神剧。
  \end{intention}

  取论域为“所有人”。记
  \begin{align*}
    &M(x): x \text{ 喜欢美剧}, \\
    &K(x): x \text{ 喜欢韩剧}, \\
    &A(x): x \text{ 喜欢抗日神剧}.
  \end{align*}
  则三条前提与结论分别形式化为
  \[
    \forall x\; (M(x) \lor K(x)),
    \qquad
    \forall x\; (A(x) \to \lnot M(x)),
    \qquad
    \exists x\; \lnot K(x),
  \]
  以及
  \[
    \exists x\; \lnot A(x).
  \]

  该结论\textbf{有效}。证明如下:

  \begin{logicproof}{4}
    \forall x\; (M(x) \lor K(x)) & premise \\
    \forall x\; (A(x) \to \lnot M(x)) & premise \\
    \exists x\; \lnot K(x) & premise \\
    \begin{subproof}
      x_0 & fresh \\
      \lnot K(x_0) & assumption \\
      M(x_0) \lor K(x_0) & \forallxelim 1 \\
      \begin{subproof}
        M(x_0) & assumption \\
        \begin{subproof}
          A(x_0) & assumption \\
          A(x_0) \to \lnot M(x_0) & \forallxelim 2 \\
          \lnot M(x_0) & \toelim 10, 9 \\
          \bot & \lnotelim 8, 11
        \end{subproof}
        \lnot A(x_0) & \lnotintro 9--12
      \end{subproof} \\
      \begin{subproof}
        K(x_0) & assumption \\
        \bot & \lnotelim 6, 14 \\
        \lnot A(x_0) & \botelim 15
      \end{subproof} \\
      \lnot A(x_0) & \lorelim 7, 8--13, 14--16 \\
      \exists x\; \lnot A(x) & \existsxintro 17
    \end{subproof} \\
    \exists x\; \lnot A(x) & \existsxelim 3, 4--18
  \end{logicproof}

  \begin{remark}
    \textbf{重点.} 由前提 (3) 先找到一个具体的人 $x_0$, 再围绕该见证人做自然推理。
    \textbf{难点.} 证明 $\lnot A(x_0)$ 时, 需要在“$M(x_0)$”这一分支内再作一次反证。
    \textbf{易错点.} 在 $\exists$-消去里引入的 $x_0$ 必须是 fresh, 不能泄漏到盒子外面。
  \end{remark}

  \textbf{分值分配（10分）.}
  \begin{enumerate}[(a)]
    \item 正确完成谓词化: \score{2}
    \item 由 $\exists x\; \lnot K(x)$ 引入 fresh 见证人: \score{2}
    \item 对 $M(x_0) \lor K(x_0)$ 正确使用 $\lor$-消去: \score{3}
    \item 在第一分支中由反证得到 $\lnot A(x_0)$: \score{2}
    \item 用 $\exists$-引入与 $\exists$-消去得到结论: \score{1}
  \end{enumerate}
\end{solution}
%%%%%%%%%%%%%%%

%%%%%%%%%%%%%%%
\newpage
\begin{problem}[一阶谓词逻辑: 自然推理系统]
  请使用一阶谓词逻辑的自然推理系统证明:
  \begin{enumerate}[(1)]
    \item $\forall x.\; (P(x) \to Q(x)) \vdash
      (\forall x.\; \lnot Q(x)) \to (\forall x\; \lnot P(x))$
    \item $\forall x\; (\phi \land \psi) \vdash
      (\forall x\; \phi) \land (\forall x\; \psi)$
    \item $(\exists x\; \phi) \lor (\exists x\; \psi) \vdash
      \exists x\; (\phi \lor \psi)$
    \item $\forall x\; \forall y\; \phi \vdash \forall y\; \forall x\; \phi$
    \item $\exists x\; \exists y\; \phi \vdash
      \exists y\; \exists x\; \phi$
  \end{enumerate}
\end{problem}

\begin{solution}
  \begin{enumerate}[(1)]
    \item
      \begin{intention}
        先假设 $\forall x\; \lnot Q(x)$, 再对任意 fresh 元 $x_0$ 证明 $\lnot P(x_0)$, 最后做两次全称/蕴含引入。
      \end{intention}

      \begin{logicproof}{4}
        \forall x\; (P(x) \to Q(x)) & premise \\
        \begin{subproof}
          \forall x\; \lnot Q(x) & assumption \\
          \begin{subproof}
            x_0 & fresh \\
            \begin{subproof}
              P(x_0) & assumption \\
              P(x_0) \to Q(x_0) & \forallxelim 1 \\
              Q(x_0) & \toelim 5, 6 \\
              \lnot Q(x_0) & \forallxelim 2 \\
              \bot & \lnotelim 8, 7
            \end{subproof}
            \lnot P(x_0) & \lnotintro 4--8
          \end{subproof} \\
          \forall x\; \lnot P(x) & \forallxintro 3--9
        \end{subproof} \\
        (\forall x\; \lnot Q(x)) \to (\forall x\; \lnot P(x)) & \tointro 2--10
      \end{logicproof}

      \begin{remark}
        \textbf{重点.} 目标是一个蕴含式, 先用 $\to$-引入; 目标内部又是全称式, 再对 arbitrary 的 $x_0$ 证明。
        \textbf{易错点.} $\forall x\; \lnot Q(x)$ 不能直接改写成 $\lnot Q(x_0)$; 必须通过 $\forall$-消去得到。
      \end{remark}

      \textbf{分值分配（10分）.}
      \begin{enumerate}[(a)]
        \item 正确使用 $\to$-引入假设 $\forall x\; \lnot Q(x)$: \score{2}
        \item 引入 fresh 的 $x_0$: \score{2}
        \item 由前提推出 $Q(x_0)$: \score{2}
        \item 与 $\lnot Q(x_0)$ 构成矛盾并得 $\lnot P(x_0)$: \score{3}
        \item 完成 $\forall$-引入与 $\to$-引入: \score{1}
      \end{enumerate}

    \item
      \begin{intention}
        从 $\forall x\; (\phi \land \psi)$ 出发, 对同一个 arbitrary 元同时抽出 $\phi$ 与 $\psi$, 然后分别做全称引入再合取。
      \end{intention}

      \begin{logicproof}{3}
        \forall x\; (\phi \land \psi) & premise \\
        \begin{subproof}
          x_0 & fresh \\
          (\phi \land \psi)[x_0/x] & \forallxelim 1 \\
          \phi[x_0/x] & \landelimleft 3 \\
          \psi[x_0/x] & \landelimright 3
        \end{subproof} \\
        \forall x\; \phi & \forallxintro 2--4 \\
        \forall x\; \psi & \forallxintro 2--5 \\
        (\forall x\; \phi) \land (\forall x\; \psi) & \landintro 7, 8
      \end{logicproof}

      \begin{remark}
        \textbf{重点.} 同一个 fresh 元 $x_0$ 上同时得到两条结论, 再分别全称化。
        \textbf{易错点.} 不能直接从 $\forall x\; (\phi \land \psi)$ 跳到 $(\forall x\; \phi) \land (\forall x\; \psi)$, 中间必须经过 $\forall$-消去和 $\forall$-引入。
      \end{remark}

      \textbf{分值分配（10分）.}
      \begin{enumerate}[(a)]
        \item 引入 fresh 的 $x_0$: \score{2}
        \item 对前提作 $\forall$-消去: \score{2}
        \item 正确抽出 $\phi[x_0/x]$: \score{2}
        \item 正确抽出 $\psi[x_0/x]$: \score{2}
        \item 完成两次 $\forall$-引入与一次合取引入: \score{2}
      \end{enumerate}

    \item
      \begin{intention}
        先对最外层析取分情况; 每个分支里再对存在量词取 fresh 见证元, 最后用同一个见证元引入 $\exists x\; (\phi \lor \psi)$。
      \end{intention}

      \begin{logicproof}{4}
        (\exists x\; \phi) \lor (\exists x\; \psi) & premise \\
        \begin{subproof}
          \exists x\; \phi & assumption \\
          \begin{subproof}
            x_0 & fresh \\
            \phi[x_0/x] & assumption \\
            \phi[x_0/x] \lor \psi[x_0/x] & \lorintroleft 4 \\
            \exists x\; (\phi \lor \psi) & \existsxintro 5
          \end{subproof} \\
          \exists x\; (\phi \lor \psi) & \existsxelim 2, 3--6
        \end{subproof} \\
        \begin{subproof}
          \exists x\; \psi & assumption \\
          \begin{subproof}
            y_0 & fresh \\
            \psi[y_0/x] & assumption \\
            \phi[y_0/x] \lor \psi[y_0/x] & \lorintroright 11 \\
            \exists x\; (\phi \lor \psi) & \existsxintro 12
          \end{subproof} \\
          \exists x\; (\phi \lor \psi) & \existsxelim 9, 10--13
        \end{subproof} \\
        \exists x\; (\phi \lor \psi) & \lorelim 1, 2--7, 9--14
      \end{logicproof}

      \begin{remark}
        \textbf{重点.} 证明顺序是“先 $\lor$-消去, 后 $\exists$-消去”。
        \textbf{难点.} 两个存在分支各自的见证元都必须 fresh, 且只在各自盒子里使用。
        \textbf{易错点.} 不能直接从 $\exists x\; \phi$ 写出 $\phi[t/x]$; 必须通过 $\exists$-消去引入新见证元。
      \end{remark}

      \textbf{分值分配（10分）.}
      \begin{enumerate}[(a)]
        \item 对最外层析取作正确分情况: \score{2}
        \item 在左分支中完成 $\exists$-消去与 $\exists$-引入: \score{3}
        \item 在右分支中完成 $\exists$-消去与 $\exists$-引入: \score{3}
        \item 两个分支都正确构造 $\phi \lor \psi$: \score{1}
        \item 用 $\lorelim$ 合并结论: \score{1}
      \end{enumerate}

    \item
      \begin{intention}
        先固定一个 arbitrary 的 $y_0$, 再固定一个 arbitrary 的 $x_0$, 由两次全称消去得到实例式, 最后反向做两次全称引入。
      \end{intention}

      \begin{logicproof}{4}
        \forall x\; \forall y\; \phi & premise \\
        \begin{subproof}
          y_0 & fresh \\
          \begin{subproof}
            x_0 & fresh \\
            \forall y\; \phi[x_0/x] & \forallxelim 1 \\
            \phi[x_0/x, y_0/y] & \forallyelim 4
          \end{subproof} \\
          \forall x\; \phi[y_0/y] & \forallxintro 3--5
        \end{subproof} \\
        \forall y\; \forall x\; \phi & \forallyintro 2--6
      \end{logicproof}

      \begin{remark}
        \textbf{重点.} 交换两个全称量词时, 只需对两个 arbitrary 元依次实例化、再依次全称化。
        \textbf{易错点.} 做 $\forall y$-引入前, $y_0$ 不能出现在任何未解除的特殊假设中。
      \end{remark}

      \textbf{分值分配（10分）.}
      \begin{enumerate}[(a)]
        \item 正确引入 fresh 的 $y_0$: \score{2}
        \item 正确引入 fresh 的 $x_0$: \score{2}
        \item 两次 $\forall$-消去得到实例式: \score{2}
        \item 完成 $\forall x$-引入: \score{2}
        \item 完成 $\forall y$-引入: \score{2}
      \end{enumerate}

    \item
      \begin{intention}
        对两个存在量词依次做消去, 取得两个 fresh 见证元后, 按相反顺序重新引入存在量词。
      \end{intention}

      \begin{logicproof}{4}
        \exists x\; \exists y\; \phi & premise \\
        \begin{subproof}
          x_0 & fresh \\
          \exists y\; \phi[x_0/x] & assumption \\
          \begin{subproof}
            y_0 & fresh \\
            \phi[x_0/x, y_0/y] & assumption \\
            \exists x\; \phi[y_0/y] & \existsxintro 5 \\
            \exists y\; \exists x\; \phi & \existsyintro 6
          \end{subproof} \\
          \exists y\; \exists x\; \phi & \existsyelim 3, 4--7
        \end{subproof} \\
        \exists y\; \exists x\; \phi & \existsxelim 1, 2--8
      \end{logicproof}

      \begin{remark}
        \textbf{重点.} 先把两个存在见证元“取出来”, 再按目标所要求的顺序“放回去”。
        \textbf{难点.} 内层 $y_0$ 的盒子结束后, 只能保留与 $y_0$ 无自由出现的结论。
        \textbf{易错点.} $\exists$-消去中的 fresh 见证元不能出现在盒子外的未解除假设与目标式中。
      \end{remark}

      \textbf{分值分配（10分）.}
      \begin{enumerate}[(a)]
        \item 正确对最外层 $\exists x$ 作消去: \score{2}
        \item 正确对内层 $\exists y$ 作消去: \score{2}
        \item 从实例式引入 $\exists x\; \phi[y_0/y]$: \score{2}
        \item 进一步引入 $\exists y\; \exists x\; \phi$: \score{2}
        \item 完成两次存在消去的封闭: \score{2}
      \end{enumerate}
  \end{enumerate}
\end{solution}
%%%%%%%%%%%%%%%

%%%%%%%%%%%%%%%%%%%%
% 如果没有需要订正的题目，可以把这部分删掉

\begincorrection
%%%%%%%%%%%%%%%%%%%%

%%%%%%%%%%%%%%%%%%%%
% 如果没有反馈，可以把这部分删掉
\beginfb

你可以写 (也可以发邮件)
\begin{itemize}
  \item 对课程及教师的建议与意见
  \item 教材中不理解的内容
  \item 希望深入了解的内容
  \item $\cdots$
\end{itemize}
%%%%%%%%%%%%%%%%%%%%
\end{document}